% =====================================================================
%  CARNET D'ENTRAINEMENT — FICHE 6 — THÈME 3
%  Niveau FACILE — Exercices
% =====================================================================
\documentclass[11pt,a4paper]{article}
\usepackage[utf8]{inputenc}
\usepackage[T1]{fontenc}
\usepackage{lmodern}
\usepackage[french]{babel}
\usepackage{amsmath,amssymb,mathtools}
\usepackage{icomma}
\usepackage[version=4]{mhchem}
\usepackage{siunitx}
\sisetup{output-decimal-marker={,}}
\usepackage{xcolor}
\usepackage{colortbl}
\usepackage{tikz}
\usepackage[most]{tcolorbox}
\usepackage{enumitem}
\usepackage{array}
\usepackage{booktabs}
\usepackage{fancyhdr}
\usepackage[hidelinks]{hyperref}
\usetikzlibrary{arrows.meta,positioning,calc,shapes.geometric}
\usepackage[a4paper,margin=2.1cm,headheight=26pt,headsep=8pt]{geometry}

\definecolor{primary}{RGB}{146,95,160}
\definecolor{primarydark}{RGB}{92,52,104}
\definecolor{excol}{RGB}{0,135,90}
\definecolor{atomO}{RGB}{220,55,45}
\definecolor{atomN}{RGB}{48,80,200}
\definecolor{atomH}{RGB}{235,235,235}
\definecolor{lonepair}{RGB}{120,94,200}
\definecolor{bondcol}{RGB}{60,60,70}

\tcbset{boxbase/.style={enhanced,breakable,boxrule=0.9pt,arc=2.5pt,
   left=3mm,right=3mm,top=2.2mm,bottom=2.2mm,
   fonttitle=\bfseries,coltitle=white,toptitle=0.6mm,bottomtitle=0.6mm}}
\newtcolorbox[auto counter]{exo}[1][]{boxbase,colback=primary!5,colframe=primary,
   title={\textsc{Exercice}~\thetcbcounter\ifx\relax#1\relax\relax\else\textmd{ —\itshape\ #1}\fi}}

\pagestyle{fancy}\fancyhf{}
\fancyhead[L]{\small\textcolor{primarydark}{\textbf{Fiche 6 · Thème 3} — Charge formelle et choix de structure}}
\fancyhead[R]{\small\textcolor{primarydark}{\textsc{Facile}}}
\fancyfoot[C]{\small\thepage}
\renewcommand{\headrulewidth}{0.8pt}
\renewcommand{\headrule}{\hbox to\headwidth{\color{primary}\leaders\hrule height \headrulewidth\hfill}}
\setlength{\parindent}{0pt}\setlength{\parskip}{3pt}

\begin{document}
\begin{center}
{\Large\bfseries\color{primarydark}Thème 3 — Niveau Facile}\\[2pt]
{\small\itshape Lire la fiche \texttt{Tuto\_Theme-3} avant de commencer. Rappel :
$q_{\mathrm{f}} = N_v - 2\times(\text{doublets libres}) - (\text{nb liaisons})$.}
\end{center}
\medskip

\begin{exo}[appliquer la formule]
Pour chaque atome, on donne le triplet $\big(N_v\ ;\ \text{nb de doublets non liants}\ ;\
\text{nb de liaisons}\big)$. Calculer sa charge formelle
$q_{\mathrm{f}} = N_v - 2\times(\text{doublets libres}) - (\text{nb liaisons})$.
\begin{enumerate}[label=\alph*),leftmargin=2em,itemsep=1pt]
  \item un \ce{O} doublement lié dans \ce{CO3^2-} : $(6\ ;\ 2\ ;\ 2)$ ;
  \item un \ce{O} simplement lié et chargé : $(6\ ;\ 3\ ;\ 1)$ ;
  \item le carbone central : $(4\ ;\ 0\ ;\ 4)$ ;
  \item l'azote de l'ion ammonium \ce{NH4+} : $(5\ ;\ 0\ ;\ 4)$.
\end{enumerate}
\end{exo}

\begin{exo}[la charge formelle de l'hydrogène]
Dans n'importe quelle liaison \ce{X-H}, l'hydrogène possède $N_v = 1$, aucun doublet non liant et
une seule liaison.
\begin{enumerate}[label=\alph*),leftmargin=2em,itemsep=1pt]
  \item Calculer la charge formelle d'un tel hydrogène « normal ».
  \item En déduire ce que vaut la charge formelle de \emph{tout} \ce{H} engagé dans une liaison simple.
\end{enumerate}
\end{exo}

\begin{exo}[la somme des charges formelles donne la charge]
On donne les charges formelles de chaque atome d'un édifice. En faisant la somme, retrouver la
\textbf{charge globale} de l'édifice.
\begin{enumerate}[label=\alph*),leftmargin=2em,itemsep=1pt]
  \item \ce{NH4+} : $q_{\mathrm{f}}(\ce{N}) = +1$ et $q_{\mathrm{f}}(\ce{H}) = 0$ pour chacun des 4 \ce{H} ;
  \item \ce{CO3^2-} : $q_{\mathrm{f}}(\ce{C}) = 0$, un \ce{O} à $0$ et deux \ce{O} à $-1$ ;
  \item \ce{OH^-} : $q_{\mathrm{f}}(\ce{O}) = -1$ et $q_{\mathrm{f}}(\ce{H}) = 0$.
\end{enumerate}
\end{exo}

\begin{exo}[la charge formelle de l'oxygène selon sa situation]
L'oxygène a toujours $N_v = 6$. Calculer sa charge formelle dans chacune de ces trois situations :
\begin{enumerate}[label=\alph*),leftmargin=2em,itemsep=1pt]
  \item \ce{O} avec \textbf{2 liaisons} et \textbf{2 doublets non liants} (oxygène « normal », ex. \ce{H2O}) ;
  \item \ce{O} avec \textbf{1 liaison} et \textbf{3 doublets non liants} ;
  \item \ce{O} avec \textbf{3 liaisons} et \textbf{1 doublet non liant} (comme dans l'ion \ce{H3O+}).
\end{enumerate}
\end{exo}

\begin{exo}[vrai ou faux ?]
Répondre par \textbf{vrai} ou \textbf{faux} et justifier en une phrase.
\begin{enumerate}[label=\alph*),leftmargin=2em,itemsep=1pt]
  \item « La charge formelle est une charge réelle portée par l'atome. »
  \item « La somme des charges formelles d'un édifice est égale à sa charge globale. »
  \item « La meilleure structure de Lewis est celle dont les charges formelles sont les plus proches de zéro. »
\end{enumerate}
\end{exo}

\end{document}
